\documentclass[12pt,a4paper]{article}
\usepackage[utf8]{inputenc}
\usepackage[spanish]{babel}
\usepackage{geometry}
\usepackage{listings}
\usepackage{xcolor}
\usepackage{parskip}
\usepackage{amsmath}
\usepackage{hyperref}

\geometry{margin=2.5cm}

\lstset{
  basicstyle=\ttfamily\small,
  keywordstyle=\color{blue}\bfseries,
  commentstyle=\color{gray},
  breaklines=true, frame=single, language=Java,
  numbers=left, numberstyle=\tiny\color{gray}
}

\title{\textbf{N-Reinas con Backtracking}\\
  \large Documentación Técnica}
\author{Inteligencia Artificial}
\date{}

\begin{document}
\maketitle
\tableofcontents
\newpage

\section{Descripción del Problema}
El problema de las \textbf{N-Reinas} consiste en colocar $N$ reinas en un tablero de $N \times N$ de forma que ninguna reina ataque a otra. Dos reinas se atacan si comparten fila, columna o diagonal.

\subsection{Número de soluciones}
\begin{center}
\begin{tabular}{|c|r|}
\hline
\textbf{N} & \textbf{Soluciones} \\
\hline
4  & 2 \\
5  & 10 \\
6  & 4 \\
7  & 40 \\
8  & 92 \\
11 & 2,680 \\
14 & 365,596 \\
\hline
\end{tabular}
\end{center}

\section{Algoritmo de Backtracking}

\subsection{Estrategia}
Se coloca una reina por fila. Para cada fila se prueba cada columna; si la posición es segura, se avanza a la siguiente fila. Si ninguna columna es válida, se retrocede (backtrack) a la fila anterior.

\subsection{Implementación}
\begin{lstlisting}
void colocarReina(int fila) {
    if (fila == N) {
        soluciones++;
        guardarSolucion(tablero.clone());
        return;
    }
    for (int col = 0; col < N; col++) {
        if (esSeguro(fila, col)) {
            tablero[fila] = col; // reina en (fila,col)
            colocarReina(fila + 1);
            tablero[fila] = -1; // backtrack
        }
    }
}
\end{lstlisting}

\subsection{Verificación de seguridad}
Una posición $(fila, col)$ es segura si ninguna reina ya colocada la ataca:

\begin{lstlisting}
boolean esSeguro(int fila, int col) {
    for (int f = 0; f < fila; f++) {
        int c = tablero[f];
        // misma columna
        if (c == col) return false;
        // diagonal (diferencia de fila == diferencia de col)
        if (Math.abs(f - fila) == Math.abs(c - col))
            return false;
    }
    return true;
}
\end{lstlisting}

No es necesario verificar filas porque se coloca exactamente una reina por fila.

\section{Visualización en Tiempo Real}

El visualizador muestra cada paso del backtracking:
\begin{itemize}
  \item \textbf{Colocación}: se dibuja la reina en color verde.
  \item \textbf{Backtrack}: la reina se elimina y se muestra en rojo momentáneamente.
  \item \textbf{Solución encontrada}: todo el tablero se resalta en amarillo.
\end{itemize}

La animación usa un \texttt{javax.swing.Timer} con delay configurable (10ms a 500ms) para controlar la velocidad de visualización.

\section{Representación del Tablero}

El tablero se representa como un arreglo \texttt{int[N]} donde \texttt{tablero[fila] = columna} indica que la reina de esa fila está en esa columna. Esto garantiza exactamente una reina por fila desde el diseño de la estructura de datos.

\section{Complejidad}

\begin{itemize}
  \item \textbf{Peor caso sin poda}: $O(N^N)$
  \item \textbf{Con backtracking}: se poda enormemente el árbol. Para $N=8$ solo se exploran $\sim$15,720 nodos de los $8^8 = 16{,}777{,}216$ posibles.
  \item \textbf{Optimización}: se pueden usar conjuntos de bits para verificar columnas y diagonales en $O(1)$.
\end{itemize}

\section{Interfaz Gráfica}

\begin{itemize}
  \item El tablero se dibuja con \texttt{paintComponent()}, alternando colores claro/oscuro.
  \item Las reinas se representan con el símbolo \texttt{Q} o con iconos personalizados.
  \item Selector de tamaño $N$ (de 4 a 14).
  \item Contador de soluciones encontradas en tiempo real.
\end{itemize}

\end{document}
